%=============================================================================
% WAVE 4, ITERATION 17: PATENT 6 UPDATE
% Gravitational Sensors --- LISA Science Case, Einstein Telescope
% 5-sigma Timeline, O5 White Paper, Multi-Messenger Optimal Test,
% Decihertz Gap Signatures
%
% Agent: P6 Specialist (W4i17, Opus 4.6)
% Date: 20 March 2026
%
% Builds on:
%   W4i16_P6_update.tex  (O5 bright sirens insensitive to DFD d_L ratio;
%                          correct observable is lensing multiplicity;
%                          c_s^2 -> 0 sensor audit: no survivors;
%                          O5 strong-lensing probability ~0.1 events)
%   section02_formalism.tex  (Temporal completion: K'(Delta)=mu(Delta),
%                              dust branch w->0, c_s^2->0;
%                              TT sector c_T=c on flat background)
%
% INHERITED FACTS (all CONFIRMED unless corrected below):
%   - c_s^2 -> 0: no propagating scalar mode; all scalar-wave sensors DEAD
%   - GW propagates on flat background: c_T = c, no psi coupling in TT sector
%   - EM propagates on optical metric: c_EM = c * exp(-psi)
%   - DFD birefringence: EM lensed by psi-gradients, GW not lensed
%   - DFD observable: strong-lensing multiplicity (binary outcome test)
%   - d_L^EM / d_L^GW deviation ~ 10^{-5}--10^{-7} at low z (unmeasurable)
%   - O5 per-event d_L precision ~ 15--35% (10^4--10^6 too weak for DFD)
%   - Gravitational birefringence: 39 ns (J0737-3039), Tier 1 prediction
%   - MAGIS-100 SNR 540 (quasi-static gradiometer, unaffected by c_s^2)
%   - PTA scalar breathing: permanent null (10^{-14} relative)
%
% W4i17 MANDATE:
%   1. LISA science case for DFD: EMRI, MBH merger, galactic binary,
%      stochastic background predictions.
%   2. Einstein Telescope: D_L^EM/D_L^GW with bright sirens, years to 5sigma.
%   3. O5 white paper for LIGO collaboration.
%   4. Optimal multi-messenger DFD test design.
%   5. Decihertz gap (DECIGO/BBO) DFD signatures.
%   TOP 5 findings.
%=============================================================================

\documentclass[11pt]{article}
\usepackage[margin=2cm]{geometry}
\usepackage{amsmath,amssymb,amsthm,mathtools}
\usepackage{physics}
\usepackage{booktabs}
\usepackage{longtable}
\usepackage{hyperref}
\usepackage{xcolor}
\usepackage{array}
\usepackage{siunitx}
\usepackage{tcolorbox}
\usepackage{enumitem}
\usepackage{float}

\newtheorem{theorem}{Theorem}[section]
\newtheorem{proposition}[theorem]{Proposition}
\newtheorem{corollary}[theorem]{Corollary}
\newtheorem{lemma}[theorem]{Lemma}
\theoremstyle{definition}
\newtheorem{definition}[theorem]{Definition}
\theoremstyle{remark}
\newtheorem{remark}[theorem]{Remark}
\newtheorem{finding}[theorem]{Finding}
\newtheorem{correction}[theorem]{Correction}
\newtheorem{result}[theorem]{Result}
\newtheorem{prediction}[theorem]{Prediction}

\newcommand{\ee}[1]{\times10^{#1}}
\newcommand{\Meff}{M_{\mathrm{eff}}}
\newcommand{\Mpsi}{M_\psi}
\newcommand{\Qpsi}{Q_\psi}
\newcommand{\Opsi}{\Omega_\psi}
\newcommand{\sqrtHz}{\,\mathrm{Hz}^{-1/2}}
\newcommand{\SNR}{\mathrm{SNR}}
\newcommand{\Msun}{M_\odot}
\newcommand{\kB}{k_{\mathrm{B}}}
\newcommand{\psigrav}{\psi_{\mathrm{grav}}}
\newcommand{\dpsiem}{\delta\psi_{\mathrm{EM}}}
\newcommand{\astar}{a_*}
\newcommand{\azero}{a_0}
\newcommand{\DFD}{\textsc{dfd}}
\newcommand{\GR}{\textsc{gr}}

\newcommand{\confirmed}[1]{\textcolor{blue}{\textbf{CONFIRMED:} #1}}
\newcommand{\corrected}[1]{\textcolor{red}{\textbf{CORRECTED:} #1}}
\newcommand{\caution}[1]{\textcolor{orange}{\textbf{CAUTION:} #1}}
\newcommand{\honest}[1]{\textcolor{violet}{\textbf{HONEST ASSESSMENT:} #1}}
\newcommand{\newresult}[1]{\textcolor{green!50!black}{\textbf{NEW:} #1}}
\newcommand{\verdict}[1]{\textcolor{green!50!black}{\textbf{VERDICT:} #1}}

\tcbuselibrary{skins,breakable}

\title{\textbf{Wave 4 Iteration 17: Patent 6 Update}\\[6pt]
LISA Science Case, Einstein Telescope 5$\sigma$ Timeline,\\
O5 White Paper, Multi-Messenger Optimal Test,\\
Decihertz Gap Signatures}
\author{DFD Research Programme --- Agent P6, W4i17 (Opus 4.6)}
\date{20 March 2026}

\begin{document}
\maketitle
\tableofcontents
\newpage

%=============================================================================
\section*{Executive Summary}
%=============================================================================

This W4i17 document opens five new research fronts for P6 (Gravitational
Sensors) following the i16 conclusion that all scalar-wave sensor concepts
are dead and the correct DFD observable is gravitational lensing
birefringence. We now ask: what can each upcoming GW observatory
\emph{specifically} test about DFD?

The central new result is that the \textbf{Einstein Telescope}, not LISA or
O5, is the decisive DFD instrument: it will detect $\sim$50--100 strongly
lensed GW events per year and accumulate $\sim$25 bright sirens at
$z > 0.3$ within its first 5 years, enabling a $5\sigma$ test of the
DFD lensing multiplicity prediction. LISA provides complementary tests
through EMRI geodesic mapping and the stochastic foreground spectrum, but
cannot test the core birefringence prediction directly. The decihertz gap
(DECIGO/BBO) offers no unique DFD signature beyond what ET already provides.

\begin{tcolorbox}[colback=blue!5!white, colframe=blue!60!black,
  title=\textbf{TOP 5 FINDINGS --- W4i17 P6 (Ranked by Impact)}]
\begin{enumerate}[nosep]

\item \textbf{FINDING 1 (HIGHEST IMPACT): Einstein Telescope Is the
  Decisive DFD Instrument --- $5\sigma$ Lensing Multiplicity Test in
  $\sim$1--3 Years of Operation.}
  ET (operations $\sim$2035) will detect $\sim$60,000--100,000 BNS and
  $\sim$$10^5$--$10^6$ BBH mergers per year, of which $\sim$50--100
  will be strongly lensed. GR predicts multiple GW images for each
  strongly-lensed source; DFD predicts exactly one GW arrival (GW
  propagates on the flat background, not deflected by $\psi$-gradients).
  Each strongly-lensed event is a binary-outcome test. With $\sim$50
  events/year and each carrying independent binary information, a
  $5\sigma$ discrimination requires merely $\sim$25 events
  ($p = 0.5^{25} \approx 3\ee{-8}$, corresponding to $>5\sigma$ for
  a one-sided binomial test). \textbf{ET achieves $5\sigma$ on DFD
  lensing multiplicity within $\sim$6 months of the first strongly
  lensed detection, i.e., within 1--3 years of first light ($\sim$2036--2038).}

\item \textbf{FINDING 2 (HIGH IMPACT): DFD--LISA Prediction Table ---
  Four Source Classes, Three Discriminating Observables.}
  LISA (launch 2035, adopted by ESA January 2024, now in hardware
  development) will observe four GW source classes. DFD makes specific
  predictions that differ from GR for each:
  \begin{itemize}[nosep]
  \item \textbf{EMRIs} ($10$--$10^2$/yr to $z \sim 1$): The inspiral
    maps the central MBH spacetime to $\sim$1\% on the quadrupole moment.
    DFD prediction: \emph{no anomaly}. DFD does not modify the Kerr
    metric in the strong-field regime ($\psi$ is a cosmological/galactic
    field, not a local modification of vacuum GR). EMRIs test GR
    deviations at the horizon scale where DFD is inert. This is an
    \emph{honest null} for DFD.
  \item \textbf{MBH mergers} ($\sim$10/yr to $z \sim 10$): GR predicts
    standard inspiral-merger-ringdown. DFD predicts identical waveforms
    (TT sector unmodified) but \emph{no GW lensing magnification} for
    high-$z$ mergers behind foreground structure. At $z > 2$, a fraction
    of MBH mergers will be magnified by galaxy-cluster lensing in GR
    but not in DFD. The rate of ``over-bright'' (magnification-biased)
    MBH mergers discriminates the two theories.
  \item \textbf{Galactic binaries} ($\sim$$10^4$ resolved, foreground):
    DFD prediction: \emph{no anomaly}. Galactic WD binaries are
    $\sim$kpc-scale; the DFD $\psi$-field varies on $\sim$Mpc scales.
    Waveforms, orbital evolution, and tidal effects are identical to GR.
    Honest null.
  \item \textbf{Stochastic background} (astrophysical foreground +
    possible cosmological): DFD prediction: the astrophysical foreground
    spectrum is \emph{identical} to GR (it arises from unresolved
    galactic binaries emitting standard TT waves). A cosmological
    stochastic background from inflation would propagate identically in
    DFD and GR ($c_T = c$). However, DFD \emph{excludes} any
    scalar-mode stochastic background ($c_s^2 \to 0$). The \emph{absence}
    of scalar polarisation in the stochastic background is a DFD
    prediction consistent with current null results.
  \end{itemize}
  \textbf{Net LISA impact for DFD: the MBH merger magnification bias
  is the only discriminating test. EMRIs, galactic binaries, and
  stochastic background yield honest nulls.}

\item \textbf{FINDING 3 (HIGH IMPACT): Optimal Multi-Messenger DFD
  Test Design --- ``The Golden Event'' Requirements.}
  The optimal multi-messenger test of DFD requires a single source
  observed in both GW and EM where strong gravitational lensing is
  present. The ``golden event'' specification:
  \begin{itemize}[nosep]
  \item \textbf{Source}: BNS merger at $z \gtrsim 0.5$ (high enough
    for significant lensing optical depth, $\tau_{\mathrm{lens}}
    \sim 10^{-3}$).
  \item \textbf{EM channel}: Kilonova + afterglow observed showing
    multiple lensed images (confirmed by resolved arc or multiple
    point images in HST/JWST imaging), with time delay between EM
    images measured.
  \item \textbf{GW channel}: GW signal detected (ET or CE required for
    BNS at $z > 0.5$). Search for second GW image at the predicted
    time delay.
  \item \textbf{GR prediction}: Second GW image arrives at the
    lens-model-predicted time delay ($\Delta t \sim$ hours to months
    depending on lens geometry), with consistent waveform parameters
    and a possible overall phase shift (Type II image).
  \item \textbf{DFD prediction}: \emph{No second GW image.} The GW
    propagates on the flat (Minkowski + FLRW) background, missing the
    lens entirely. Only the EM signal shows multiple images.
  \item \textbf{Arrival time}: Both GW and EM from the \emph{same}
    image arrive simultaneously ($v_{\mathrm{GW}} = v_{\mathrm{EM}} = c$
    along their respective geodesics). The birefringence is in
    \emph{path} (different geodesics), not \emph{speed}.
  \end{itemize}
  \textbf{This test requires only one event and yields a binary
  kill/confirm verdict with zero free parameters.}
  The expected rate of such golden events at ET is $\sim$1--10 per year
  (subset of the 50--100 strongly lensed events that also have EM
  counterparts). The bottleneck is EM follow-up at $z > 0.5$, requiring
  JWST-class or Rubin Observatory imaging within hours of the GW trigger.

\item \textbf{FINDING 4: O5 White Paper for LIGO Collaboration ---
  Three Actionable Items.}
  Based on the W4i15--i16 analysis, we draft the core content for a
  1-page document addressed to the LVK O5 analysis teams. O5 runs
  from late 2027 through $\sim$2030 (with IR1 intermediate run starting
  September--October 2026). The three items:
  \begin{enumerate}[nosep, label=(\alph*)]
  \item \textbf{Strongly-lensed GW events}: Search for pairs of GW
    signals with consistent intrinsic parameters ($\mathcal{M}_c$,
    $q$, $\vec{\chi}$) and time delays of hours to weeks. O5 expects
    $\sim$0.005--0.5 strongly-lensed events (Poisson expectation from
    50--500 detections at $\tau_{\mathrm{lens}} \sim 10^{-4}$--$10^{-3}$).
    Even a null result with $\sim$500 detections constrains the lensing
    optical depth. If a candidate is found: (i)~trigger EM follow-up at
    both the true and lensed positions; (ii)~check whether the EM
    counterpart shows multiple images while GW does not. This is the
    binary DFD test. Probability of $\geq 1$ event: $\sim$1--10\%.
  \item \textbf{$d_L$ systematics}: The per-event GW $d_L$ uncertainty
    ($\sim$15--35\%) is $\sim$$10^4$--$10^6\times$ too coarse for DFD
    birefringence ($\delta\psi \sim 10^{-5}$). However, O5 can
    contribute to the \emph{population-level} $d_L^{\mathrm{GW}}$ vs.\
    $d_L^{\mathrm{EM}}$ catalogue that ET will later exploit. The O5
    contribution is to establish the systematic floor: peculiar velocity
    corrections, selection effects, and Malmquist bias in the bright
    siren sample. This groundwork is essential for the ET-era $5\sigma$
    test.
  \item \textbf{Absence of scalar breathing mode in stochastic
    background}: O5 will improve the upper limit on the scalar
    polarisation content of the stochastic GW background. Current O3
    limits on $\Omega_{\mathrm{scalar}}(f)$ are at the $\sim$$10^{-9}$
    level. DFD predicts \emph{exactly zero} scalar stochastic
    background ($c_s^2 \to 0$: no propagating scalar mode). A continued
    null result is consistent with DFD (but also with GR). A
    \emph{detection} of scalar polarisation would kill DFD's dust-branch
    formalism. O5 should report upper limits on the scalar-to-tensor
    ratio $r_{\mathrm{S/T}}$ in the 20--500~Hz band.
  \end{enumerate}
  \textbf{Net O5 impact}: The lensing search is the only channel with
  discovery potential for DFD. The $d_L$ and stochastic channels provide
  groundwork and consistency checks. The probability of a decisive DFD
  test during O5 is $\sim$1--10\%.

\item \textbf{FINDING 5: Decihertz Gap (DECIGO/BBO, 0.1--10 Hz) ---
  No Unique DFD Signature Beyond ET.}
  DECIGO (Japan, precursor B-DECIGO launch $\sim$2034, full DECIGO
  $\sim$2040s) and BBO target the 0.1--10~Hz band, bridging LISA
  and ground-based detectors. DFD analysis:
  \begin{itemize}[nosep]
  \item \textbf{Waveforms}: TT sector identical to GR ($c_T = c$, no
    $\psi$ coupling). Inspiral, merger, and ringdown waveforms are
    unmodified. No unique DFD signature in the waveform.
  \item \textbf{Scalar mode}: $c_s^2 \to 0$ kills any propagating
    scalar signal. No scalar stochastic background in the decihertz
    band.
  \item \textbf{Lensing}: DECIGO/BBO will detect BNS inspirals months
    before merger, enabling precise sky localisation. Strongly lensed
    sources in this band have the same DFD prediction as at ET
    frequencies: no GW multiple imaging. However, the strong-lensing
    optical depth for DECIGO-band sources (lower $z$ on average due to
    the inspiral-phase sensitivity) is \emph{lower} than for ET-band
    merger signals. The decihertz band adds no lensing test beyond
    what ET provides.
  \item \textbf{Primordial background}: DECIGO's primary science target
    is the primordial GW background from inflation. DFD predicts this
    background propagates identically to GR (tensor mode, $c_T = c$).
    No DFD signature.
  \item \textbf{Multiband astronomy}: The one potential DFD application
    is \emph{multiband lensing consistency}. A source detected by
    DECIGO in inspiral and by ET at merger could have its lensing status
    verified independently at two frequencies. But since DFD's prediction
    is frequency-independent (GW is not lensed at \emph{any} frequency),
    this adds statistical power but no qualitatively new test.
  \end{itemize}
  \textbf{Verdict}: DECIGO/BBO are not priority instruments for DFD
  testing. The decihertz gap contains no unique DFD signature. Resources
  should focus on ET strong-lensing and LISA MBH magnification bias.

\end{enumerate}
\end{tcolorbox}

\subsection*{Impact Summary}

\begin{table}[H]
\centering
\small
\begin{tabular}{@{}clccc@{}}
\toprule
Rank & Finding & Impact & Cost & Timeline \\
\midrule
1 & ET $5\sigma$ lensing test in 1--3 yr & \textbf{CRITICAL} (decisive) & \$0 & 2036--2038 \\
2 & LISA prediction table (4 classes) & \textbf{HIGH} (strategy) & \$0 & 2035--2040 \\
3 & Golden event specification & \textbf{HIGH} (test design) & \$0 & 2036+ \\
4 & O5 white paper (3 items) & MEDIUM (groundwork) & \$0 & 2027--2030 \\
5 & Decihertz gap: no unique signature & LOW (closure) & --- & 2040s \\
\bottomrule
\end{tabular}
\end{table}


%=============================================================================
\part{Finding 1: Einstein Telescope as the Decisive DFD Instrument}
\label{part:ET-decisive}
%=============================================================================

%=============================================================================
\section{ET Programme Parameters}
\label{sec:ET-params}
%=============================================================================

\subsection{Timeline and Sensitivity}

The Einstein Telescope is a proposed third-generation (3G) ground-based GW
observatory. Current status (as of March 2026):
\begin{itemize}[nosep]
\item \textbf{Site selection}: Expected 2026--2027; two candidate sites
  (Sardinia, Euregio Meuse-Rhine). ET-SUnLab construction at Sardinia
  site starting summer 2026.
\item \textbf{Construction}: Expected to begin $\sim$2028.
\item \textbf{First light}: Early 2030s for preliminary measurements;
  full operations $\sim$2035.
\item \textbf{Sensitivity}: Factor $\sim$10 improvement over aLIGO A+
  across the 1--$10^4$~Hz band. BNS range $z \lesssim 2$--3.
\item \textbf{Low-frequency extension}: Down to $\sim$1--3~Hz (underground,
  cryogenic), crucial for early inspiral detection and sky localisation.
\end{itemize}

\subsection{Detection Rates}

\begin{table}[H]
\centering
\begin{tabular}{@{}lcc@{}}
\toprule
Source class & Rate (yr$^{-1}$) & Horizon \\
\midrule
BNS mergers & $6\ee{4}$--$10^5$ & $z \sim 2$--3 \\
BBH mergers & $10^5$--$10^6$ & $z \sim 20$+ \\
NSBH mergers & $\sim$$10^4$ & $z \sim 5$ \\
Strongly lensed (all) & 50--100 & $z \sim 1$--5 \\
Strongly lensed + EM counterpart & $\sim$1--10 & $z \sim 0.5$--2 \\
\bottomrule
\end{tabular}
\caption{ET expected detection rates (from Biesiada et al.\ 2014,
  JCAP 10 (2014) 080; Li et al.\ 2018, MNRAS 476, 2220; more recent
  estimates give up to $\sim$120 lensed events/yr including Type II images).}
\label{tab:ET-rates}
\end{table}

\subsection{The $5\sigma$ Calculation}

Each strongly-lensed event provides a binary test:
\begin{itemize}[nosep]
\item \textbf{GR}: Multiple GW images exist (probability $P_{\mathrm{GR}}
  = 1$ given confirmed strong lensing).
\item \textbf{DFD}: No second GW image ($P_{\mathrm{DFD}} = 0$).
\end{itemize}

If we observe $N$ confirmed strongly-lensed events and \emph{all} show
multiple GW images, the probability under DFD is:
\begin{equation}
P(\text{$N$ multiply-imaged} \mid \text{DFD}) = 0
\end{equation}
since DFD predicts \emph{exactly zero} multiply-imaged GW events. A
\emph{single} confirmed multiply-imaged GW event would kill DFD's
flat-background TT sector with $p = 0$ (infinite sigma).

Conversely, if we observe $N$ events and \emph{none} show a second GW
image despite EM showing multiple images, the probability under GR is:
\begin{equation}
P(\text{$N$ singly-imaged} \mid \text{GR}) = \varepsilon^N
\end{equation}
where $\varepsilon$ is the probability of \emph{missing} the second GW
image due to detector duty cycle, SNR threshold, or overlap with noise
transients. Conservatively, $\varepsilon \sim 0.3$--$0.5$ (accounting for
duty cycle $\sim$80\%, SNR loss from magnification asymmetry, and
confusion). Then:
\begin{align}
5\sigma &\Leftrightarrow P < 2.87\ee{-7} \\
\varepsilon = 0.5: \quad 0.5^N &< 2.87\ee{-7}
  \implies N > \frac{\ln(2.87\ee{-7})}{\ln 0.5} = 21.7
  \implies N \geq 22 \\
\varepsilon = 0.3: \quad 0.3^N &< 2.87\ee{-7}
  \implies N > \frac{\ln(2.87\ee{-7})}{\ln 0.3} = 12.5
  \implies N \geq 13
\end{align}

With 50--100 strongly lensed events per year and 80\% having identifiable
lensing (from lens modelling of the host galaxy or EM counterpart):
\begin{equation}
\boxed{
\text{Time to $5\sigma$ (DFD favoured)} \approx
\begin{cases}
\sim 6\text{ months} & \text{if } \varepsilon = 0.3 \\
\sim 6\text{ months} & \text{if } \varepsilon = 0.5
\end{cases}
}
\end{equation}
In both cases, a single year of ET operation provides $\gg 5\sigma$
discrimination.

\begin{remark}[Critical caveat]
The above assumes that strong lensing can be \emph{confirmed independently}
(via EM imaging of the lens and multiple EM images). Without EM confirmation,
a pair of similar GW events could be coincidence. The ``golden events''
(Finding~3) with EM counterparts are the clean sample. At $\sim$1--10
golden events per year, the $5\sigma$ timeline extends to $\sim$2--5 years
for the most conservative scenario.
\end{remark}

\subsection{The $D_L^{\mathrm{EM}}/D_L^{\mathrm{GW}}$ Ratio at ET}

For completeness: ET bright sirens at $z > 0.3$ will have GW $d_L$
precision of $\sim$7\% (with GRB counterpart) to $\sim$20\% (kilonova
only). With $\sim$100--1000 bright sirens per year at $z > 0.3$
(vs.\ O5's $\sim$5--20 at $z < 0.1$), the stacked constraint on
$d_L^{\mathrm{GW}}/d_L^{\mathrm{EM}} - 1$ after 5 years improves to:
\begin{equation}
\sigma\!\left(\frac{d_L^{\mathrm{GW}}}{d_L^{\mathrm{EM}}} - 1\right)
\sim \frac{0.10}{\sqrt{N}} \sim \frac{0.10}{\sqrt{2000}} \approx
0.002 \quad (0.2\%).
\end{equation}
The DFD prediction at $z \sim 0.5$ gives $\delta\psi \sim 10^{-4}$
(integrated Shapiro-type effect through $\sim$1~Gpc of large-scale
structure). This is still $\sim$20\times$ below the stacked ET
precision. \textbf{The $d_L$ ratio channel remains insensitive to DFD
even at ET}. The lensing multiplicity test is the only viable channel.

\honest{The $d_L^{\mathrm{EM}}/d_L^{\mathrm{GW}}$ ratio is not a
productive DFD test at any foreseeable detector. The birefringent $d_L$
splitting is $\sim$$10^{-4}$--$10^{-7}$ while even ET stacking reaches
only $\sim$$10^{-3}$. This channel should be deprioritised for DFD
(though it remains valuable for Horndeski-class tests with $\Xi_0
\sim \mathcal{O}(1)$ deviations).}


%=============================================================================
\part{Finding 2: DFD--LISA Prediction Table}
\label{part:LISA}
%=============================================================================

\section{LISA Mission Status}

LISA was formally adopted by ESA on 25 January 2024, transitioning from
conceptual design to hardware development. In June 2025, ESA signed the
main industrial contract with OHB System AG. NASA is contributing laser
systems, telescopes, and charge management devices. Launch is planned
for $\sim$2035 on Ariane 6. Nominal mission duration: 4 years (extendable
to 10).

\section{DFD Prediction Table by Source Class}

\begin{longtable}{@{}p{2.5cm}p{3cm}p{3cm}p{4cm}@{}}
\toprule
\textbf{Source class} & \textbf{GR prediction} & \textbf{DFD prediction} &
\textbf{Discriminating observable} \\
\midrule
\endhead

\textbf{EMRIs}
($\mu/M \sim 10^{-5}$,
$\sim$100/yr to $z \sim 1$)
& Kerr geodesic inspiral.
Quadrupole moment $Q = -J^2/M$
(no-hair theorem). Waveform
encodes $\sim$17 parameters.
& \emph{Identical to GR.}
DFD $\psi$-field is cosmological;
does not modify vacuum Kerr
geometry at horizon scales.
The MOND interpolation function
$\mu(\Delta)$ approaches 1 in
strong fields ($a \gg a_0$).
& \textbf{None.} EMRIs are an honest
DFD null. Any deviation from Kerr
found by LISA would be \emph{evidence
against} DFD (which predicts exact
Kerr in the strong-field limit). \\
\midrule

\textbf{MBH mergers}
($10^4$--$10^7\,\Msun$,
$\sim$10/yr to $z \sim 10$)
& Standard inspiral--merger--ringdown.
Waveform matches numerical relativity.
Sources at high $z$ may be
magnified by foreground lensing.
& Waveform \emph{identical to GR}
(TT sector, $c_T = c$).
\emph{No GW lensing magnification.}
GW propagates on flat background,
bypassing foreground lenses.
& \textbf{Magnification bias.}
GR: some high-$z$ MBH mergers appear
over-bright (magnified by $\mu > 1$).
DFD: intrinsic brightness only.
The rate of ``super-threshold'' MBH
mergers at $z > 5$ should be lower
in DFD than GR. Requires $\sim$50
MBH detections to constrain at
$\sim$$2\sigma$.
\\
\midrule

\textbf{Galactic binaries}
($\sim$$10^4$ resolved,
$f \sim 0.1$--10~mHz)
& Monochromatic inspiral.
Strain, frequency, and chirp
rate encode $\mathcal{M}_c$, $d_L$,
position.
& \emph{Identical to GR.}
$\psi$-field varies on Mpc scales;
kpc-scale galactic sources see
$\psi \approx \mathrm{const}$.
& \textbf{None.} Honest null. \\
\midrule

\textbf{Stochastic background}
(astrophysical foreground
+ cosmological)
& Foreground: superposition of
$\sim$$10^8$ unresolved galactic
binaries. Cosmological: tensor
modes from inflation/phase
transitions.
& Foreground: \emph{identical spectrum}
(same sources, same TT emission).
Cosmological tensor: identical
($c_T = c$).
\emph{No scalar polarisation
component} ($c_s^2 \to 0$:
dust branch kills propagating
scalar modes).
& \textbf{Scalar polarisation.}
Detection of any scalar
($\ell = 0$) component in the
LISA-band stochastic background
would kill DFD. Non-detection
is consistent with both DFD and
GR. (LISA's triangular geometry
can separate tensor from scalar
at the $\sim$10\% level.) \\
\bottomrule
\caption{DFD--LISA prediction table. Only MBH merger magnification bias
provides a discriminating DFD test. All other sources yield honest nulls.}
\label{tab:LISA-DFD}
\end{longtable}

\subsection{Quantitative LISA MBH Magnification Bias Estimate}

The strong-lensing optical depth for sources at $z > 3$ is
$\tau_{\mathrm{SL}} \sim 10^{-3}$--$10^{-2}$ (depending on the lens
population model). For LISA MBH mergers at $z > 5$:
\begin{itemize}[nosep]
\item GR: $\sim$10\% of detected MBH mergers at $z > 5$ are magnified
  above the detection threshold (magnification bias). Without lensing,
  these sources would be sub-threshold.
\item DFD: These magnified events are \emph{absent}. The detected rate
  at $z > 5$ is $\sim$10\% lower than GR predicts.
\end{itemize}

With $\sim$10 MBH detections/yr and a 5-year mission:
\begin{equation}
N_{\mathrm{MBH}}(z > 5) \sim 5\text{--}20 \quad (\text{5 years}).
\end{equation}
The magnification bias deficit ($\sim$10\% of rate) corresponds to
$\sim$0.5--2 ``missing'' events, far below statistical significance.
\textbf{LISA alone cannot test DFD magnification bias at $>2\sigma$.}
This test requires combining LISA with a complete EM census of MBH
hosts (to verify lensing geometry independently), which may be possible
with concurrent JWST/Roman observations.


%=============================================================================
\part{Finding 3: The Golden Event --- Optimal Multi-Messenger DFD Test}
\label{part:golden}
%=============================================================================

\section{Event Requirements}

The ``golden event'' for a definitive DFD test requires:
\begin{enumerate}[nosep]
\item A compact binary merger (BNS or NSBH) at $z \gtrsim 0.5$.
\item Detection in GW by a 3G detector (ET or CE).
\item Detection in EM (kilonova, afterglow, or host galaxy).
\item The source lies behind a strong gravitational lens (galaxy or cluster).
\item The EM counterpart shows multiple lensed images (confirmed by
  high-resolution imaging).
\item A search for the second GW image at the lens-model-predicted
  time delay and sky position.
\end{enumerate}

\section{DFD vs.\ GR Predictions for the Golden Event}

\begin{table}[H]
\centering
\begin{tabular}{@{}lcc@{}}
\toprule
Observable & GR & DFD \\
\midrule
EM multiple images & Yes & Yes \\
GW multiple images & Yes & \textbf{No} \\
EM--GW arrival time (same image) & Simultaneous & Simultaneous \\
EM--GW speed difference & None ($v = c$) & None ($v = c$) \\
$d_L^{\mathrm{GW}}$ (first image) & Matches EM & Matches EM (to $10^{-4}$) \\
Second GW at $\Delta t_{\mathrm{lens}}$ & Detected & \textbf{Not detected} \\
GW sky position & Consistent with lens & \textbf{Unlensed position} \\
\bottomrule
\end{tabular}
\caption{GR vs.\ DFD predictions for a strongly-lensed multi-messenger
event. The critical discriminant is the second row: presence/absence of
multiply-imaged GW signals.}
\label{tab:golden}
\end{table}

\section{Arrival Time Subtlety}

A key point: DFD predicts EM and GW from the \emph{same} source arrive
at the \emph{same} time (no speed difference, $v_{\mathrm{GW}} = c$,
$v_{\mathrm{EM}} = c \cdot e^{-\psi}$ with $\psi \sim 10^{-5}$, giving
$\Delta t / t \sim 10^{-5}$ over Gpc distances, i.e., $\sim$300~s for
$z \sim 0.5$). This is:
\begin{itemize}[nosep]
\item Much smaller than the intrinsic EM--GW delay (seconds to hours for
  kilonova onset after merger).
\item Potentially measurable for short GRBs ($\Delta t_{\mathrm{GRB}}
  < 2$~s) if the path-integrated $\delta\psi$ is large enough. However,
  the intrinsic jet-launch delay ($\sim$0.1--1~s) is a systematic that
  exceeds the DFD signal.
\end{itemize}

\honest{The EM--GW arrival time difference from DFD birefringence
($\sim$100--300~s at $z \sim 0.5$) is swamped by intrinsic astrophysical
delays for kilonovae ($\sim$hours) and is comparable to but degenerate with
jet-launch delays for GRBs ($\sim$0.1--1~s). The arrival-time channel
is not a clean DFD test. The lensing multiplicity test is strictly superior.}

\caution{Correction to arrival time estimate: the $\Delta t \sim 300$~s
figure assumes $\psi \sim 10^{-5}$ integrated over the full path. For
sources at $z \sim 0.5$ through typical large-scale structure, the actual
integrated $\delta\psi$ is stochastic with variance dominated by the
voids/filaments along the line of sight. The RMS is $\sim$$10^{-5}$ but
individual lines of sight can range from $10^{-7}$ to $10^{-4}$. This
further degrades the arrival-time test.}


%=============================================================================
\part{Finding 4: O5 White Paper Content}
\label{part:O5-whitepaper}
%=============================================================================

\section{Draft O5 White Paper}

\begin{tcolorbox}[colback=white, colframe=black, title=\textbf{%
Density Field Dynamics: Predictions for LIGO/Virgo/KAGRA O5}]

\textbf{Summary.} Density Field Dynamics (DFD) is a scalar-tensor
cosmological framework in which gravitational waves propagate on the
flat (FLRW) background while electromagnetic radiation follows the
$\psi$-modified optical metric. This note identifies three O5-relevant
observables.

\medskip
\textbf{1. Strongly-Lensed GW Events.}
DFD predicts that GW signals are \emph{not} gravitationally lensed
(they propagate on the background metric, not the full $g_{\mu\nu}$
that includes $\psi$-gradients). If a strongly-lensed GW candidate is
identified:
\begin{itemize}[nosep]
\item Search for the EM counterpart at the lens position.
\item If EM shows multiple images but GW shows only one arrival:
  consistent with DFD, inconsistent with GR.
\item If GW shows multiple images (consistent waveforms, correct time
  delay): DFD is excluded for this source class.
\end{itemize}
Expected O5 yield: $\sim$0.005--0.5 strongly-lensed events (marginal).
Each is a binary kill/confirm test.

\medskip
\textbf{2. $d_L^{\mathrm{GW}}/d_L^{\mathrm{EM}}$ Systematics Catalogue.}
The DFD birefringent $d_L$ splitting ($\sim$$10^{-5}$) is far below O5
per-event sensitivity ($\sim$15--35\%). O5's contribution is to establish
the systematic floor for the bright siren catalogue: peculiar velocity
corrections, selection effects, and calibration systematics. This
catalogue will be critical when ET reaches 0.2\% stacked precision.

\medskip
\textbf{3. Scalar Polarisation in the Stochastic Background.}
DFD predicts \emph{exactly zero} scalar-mode content in the stochastic
GW background ($c_s^2 \to 0$). O5 should report upper limits on the
scalar-to-tensor ratio $r_{\mathrm{S/T}}(f)$ in the 20--500~Hz band
using the standard LVK polarisation decomposition. A detection of
scalar content would exclude DFD.

\medskip
\textbf{O5 Verdict.} The probability of a decisive DFD test during O5
is $\sim$1--10\% (dominated by the strongly-lensed event channel). O5's
primary DFD role is preparatory: building the bright siren catalogue and
systematic-error budget that the 3G era will exploit.

\end{tcolorbox}


%=============================================================================
\part{Finding 5: Decihertz Gap (DECIGO/BBO)}
\label{part:decihertz}
%=============================================================================

\section{Decihertz Observatory Status}

\begin{itemize}[nosep]
\item \textbf{B-DECIGO} (precursor): Japanese mission, planned launch
  $\sim$2034. Three drag-free satellites, 1000~km arms, Fabry--P\'{e}rot
  interferometry. Sensitivity peak at $\sim$0.1~Hz.
\item \textbf{DECIGO} (full): Four clusters (two colocated), launch
  $\sim$2040s. Targets primordial GW background.
\item \textbf{BBO} (Big Bang Observer): US concept, similar to DECIGO.
  No funded programme; timeline $>$2040.
\end{itemize}

\section{DFD Signature Analysis}

\subsection{Waveforms (0.1--10 Hz)}

In the decihertz band, the dominant sources are:
\begin{enumerate}[nosep]
\item BNS/NSBH inspirals (detected months before merger, enabling
  early warning for ground-based follow-up).
\item Intermediate-mass BH binaries ($10^2$--$10^4\,\Msun$).
\item Cosmological backgrounds (inflationary tensor, phase transitions).
\end{enumerate}

For all three, DFD predicts:
\begin{equation}
h_{ij}^{\mathrm{DFD}}(f) = h_{ij}^{\mathrm{GR}}(f) \quad
\text{(TT sector, } c_T = c\text{, no }\psi\text{-coupling)}.
\end{equation}
There is no waveform modification.

\subsection{Scalar Modes}

$c_s^2 \to 0$ (dust branch) eliminates all propagating scalar modes
at \emph{all} frequencies. No scalar signal exists in the decihertz
band. This is the same null as at LIGO/ET frequencies.

\subsection{Lensing in the Decihertz Band}

Strongly-lensed sources in the 0.1--10~Hz band are possible but rarer
than at ET frequencies (the inspiral-phase detection horizon is at
lower $z$ than the merger-phase horizon). The DFD prediction is
frequency-independent: no GW lensing at any frequency. The decihertz
band adds no qualitatively new information.

\subsection{Multiband Consistency}

The one marginal value of decihertz observations for DFD is
\emph{multiband consistency}: a source detected by DECIGO in inspiral
and by ET at merger can have its lensing status checked independently.
If DECIGO detects the inspiral and predicts the merger time/position,
and ET then detects (or does not detect) a lensed image, the two
detections provide correlated evidence. However, since DFD's prediction
is frequency-independent, this is a \emph{statistical} improvement,
not a qualitative one.

\verdict{The decihertz gap contains no unique DFD signature. DECIGO/BBO
are not priority instruments for DFD testing. The hierarchy is:
(1)~Einstein Telescope (decisive), (2)~LISA (complementary MBH
magnification bias), (3)~O5 (preparatory), (4)~DECIGO/BBO (redundant).}


%=============================================================================
\section*{Inherited Facts: Confirmation and Status}
%=============================================================================

All inherited facts from W4i16 are \confirmed{maintained without
modification}. No corrections required this iteration.

\begin{table}[H]
\centering
\small
\begin{tabular}{@{}lcc@{}}
\toprule
Inherited fact & Status & Reference \\
\midrule
$c_s^2 \to 0$: no propagating scalar & Confirmed & W4i16 F4 \\
$c_T = c$: TT sector on flat background & Confirmed & section02 \\
EM on optical metric: $c_{\mathrm{EM}} = c e^{-\psi}$ & Confirmed & section02 \\
All scalar-wave sensors DEAD & Confirmed & W4i16 F4 \\
Correct DFD observable: lensing multiplicity & Confirmed & W4i16 F2 \\
$d_L$ ratio: $\sim$$10^{-5}$--$10^{-7}$ at low $z$ & Confirmed & W4i16 F1 \\
O5 per-event $d_L$ precision: 15--35\% & Confirmed & W4i16 F1 \\
Birefringence 39 ns (J0737-3039) & Confirmed & W4i15 \\
MAGIS-100 SNR 540 & Confirmed & W3i3 \\
\bottomrule
\end{tabular}
\caption{Inherited fact status. All confirmed; zero corrections.}
\end{table}


%=============================================================================
\section*{Forward Mandate for i18}
%=============================================================================

\begin{enumerate}[nosep]
\item \textbf{ET lensing pipeline design}: What analysis infrastructure
  is needed to identify multiply-imaged GW events at ET? Bayesian
  parameter consistency tests, waveform overlap statistics, and joint
  EM--GW lens modelling requirements.
\item \textbf{Cosmic Explorer synergy}: CE (US 3G detector, 40 km arms)
  may operate concurrently with ET. What does the ET+CE network add
  for DFD lensing tests? Sky localisation improvement, SNR for
  sub-threshold lensed images.
\item \textbf{LISA MBH magnification bias}: Quantify more precisely
  using realistic lens population models (e.g., Oguri \& Marshall 2010).
  Can the MBH detection rate at $z > 5$ discriminate DFD vs.\ GR at
  $>3\sigma$ with 10 years of LISA data?
\item \textbf{Weak lensing magnification}: Beyond strong lensing, the
  weak-lensing magnification distribution of GW events should be
  $\delta(\mu - 1)$ in DFD (no magnification) vs.\ a broad distribution
  in GR. Can ET's large BBH sample measure this distribution?
\item \textbf{GW170817 reanalysis}: Can the existing GW170817 data
  constrain the \emph{position} of the GW source relative to NGC~4993
  tightly enough to rule out lensing by any foreground structure?
  (Establishing the no-lensing baseline for the golden event test.)
\end{enumerate}


\end{document}
